\documentclass{article}%
\usepackage[T1]{fontenc}%
\usepackage[utf8]{inputenc}%
\usepackage{lmodern}%
\usepackage{textcomp}%
\usepackage{lastpage}%
\usepackage{geometry}%
\geometry{margin=2.2cm}%
%
\usepackage{amsmath}%
\usepackage{amssymb}%
\usepackage{booktabs}%
\usepackage{longtable}%
\usepackage{array}%
\usepackage{hyperref}%
\hypersetup{hidelinks}%
\title{ITA 2026 {-} Química (claude\_sonnet\_5\_5) x Poliedro\\[4pt]\large Comparação com o gabarito Poliedro Resolve\thanks{Gerado com gpt-resolve: \url{https://github.com/lgabs/gpt-resolve}}}%
\date{10/11/2026}%
%
\begin{document}%
\normalsize%
\maketitle%
\section*{Resumo}\textbf{Nota média: 10,0/10} (32 de 32 itens com a resolução correta; cada item vale o mesmo dentro da questão).\par\medskip%
\begin{longtable}{@{}c c p{0.78\textwidth}@{}}\toprule Questão & Nota & Resumo \\ \midrule
1 & 10,0 & Nos três itens a resolução bate com o gabarito Poliedro: a) 0,40 mol/L; b) 2,0 mol/L; c) $\approx$ 2,1$\times$10$^{-}$$^2$ mol/L (2/95 = 0,8/38). Refiz as contas e as duas estão certas. A hidrólise de A$^{-}$ (K\_h = 10$^{-}$$^2$) com c\_sal = 2$\times$10$^{-}$$^2$ mol/L dá exatamente {[}OH$^{-}${]} = 10$^{-}$$^2$ mol/L, isto é, pH 12. A resolução e o gabarito tratam corretamente essa hidrólise de 50\%, sem a aproximação {[}A$^{-}${]} $\approx$ c\_sal. Fora isso, há só erros de digitação, no enunciado (?Solução I?) e na transcrição do gabarito, sem efeito no resultado. \\
2 & 10,0 & A questão 2 bate com o gabarito do Poliedro Resolve nos dois itens. No item a, os dois chegam a 0,025 mol de O2 e V $\approx$ 0,61 L. O gabarito dá 0,6109 L porque usou R = 0,082; a resolução dá 0,6116 L porque usou R = 0,0821, que é o valor da prova. No item b, os dois usam cinética de 1ª ordem com kt = 0,3465. A resolução faz 0,5$\cdot$0,707 = 0,3535 mol$\cdot$L$^{-}$$^{1}$ e o gabarito faz 0,707$^3$ = 0,353, e as duas contas são equivalentes porque 0,707$^2$ $\approx$ 0,5. O único deslize é do gabarito, que escreveu a unidade como 'mol' em vez de mol$\cdot$L$^{-}$$^{1}$. A resolução está correta e completa nos dois itens. \\
3 & 10,0 & A resolução está correta nos dois itens e concorda com o gabarito do Poliedro. No item a, ambos dão t $\approx$ 553,6 s ($\approx$ 554 s); o valor exato, sem arredondar $\surd$5, é 552,8 s. No item b, o modelo e o equacionamento são iguais (V$_{2}$ = V$_{1}$/2, precipitação até {[}A?{]}{[}B$^{-}${]} = K\_PS). A resolução chega à forma exata {[}A?{]} = (1+$\surd$3)$\surd$5$\cdot$10$^{-}$$^3$ $\approx$ 6,1$\cdot$10$^{-}$$^3$ mol/L e {[}B$^{-}${]} = ($\surd$3-1)$\surd$5$\cdot$10$^{-}$$^3$ $\approx$ 1,6$\cdot$10$^{-}$$^3$ mol/L. O gabarito dá 6,135 e 1,655$\cdot$10$^{-}$$^3$ mol/L; a pequena diferença vem dos arredondamentos que ele faz no meio da conta, e o produto dos valores dele nem chega a fechar exatamente em K\_PS. Não há erro conceitual nem de conta na resolução. \\
4 & 10,0 & A resolução e o gabarito concordam nos dois itens. No item a, a resolução obtém 5,58 g, a partir do teor medido de 28,5\%, e o gabarito 5,57 g, a partir da fórmula com n = 9. As respostas se equivalem: a diferença de cerca de 0,25\% vem só do caminho de cálculo, e as contas estão certas nos dois. No item b, os dois dão n = 9, confirmado pela conta: n = 9,06, e o teor teórico com n = 9 é 28,57\%. O gabarito tem um erro de digitação no item b: falta o fator 2 na equação escrita. Isso não altera o resultado. A resolução está correta e completa. \\
5 & 10,0 & A resolução está correta nos quatro itens e bate com o gabarito do Poliedro. Item a: 3 arranjos, um para cada átomo central. Item b: BF$_{3}$ e NH$_{3}$, que o gabarito aceita. Item c: nitração do fluorobenzeno dando os isômeros orto e para. Item d: H$_{2}$S < CH$_{3}$OH < H$_{2}$O. Conferi o enunciado no PDF original (página 7). Também conferi as cargas formais, a orientação orto/para do F e os pontos de ebulição. Não achei erro na resolução nem no gabarito. As ressalvas são de baixa importância: a ambiguidade de 'arranjos atômicos' (contar ou não o anel) e o rótulo de 'minoritário' dado ao isômero orto. \\
6 & 10,0 & A resolução bate com o gabarito do Poliedro nos três itens. Os quatro valores derivados são iguais nos dois: pKb(CH3CC-) = -11, pKa(H3O+) = -1,70, pKa(NH4+) = 9,25 e pKa(NH3) = 36. As constantes das reações (10\^{}-26,7, 10\^{}-15,75 e 10\^{}11) e as conclusões também coincidem: a) não favorece os produtos, b) não favorece, c) favorece. Refiz as contas e todas estão certas. As duas usam o mesmo critério: o equilíbrio se desloca para o lado do ácido mais fraco, com K = 10\^{}(pKa do ácido nos produtos - pKa do ácido nos reagentes). O único ponto em aberto é o comentário do gabarito sobre o pKb da água (15,70 contra 14), que é questão de convenção e não muda nenhuma conclusão. \\
7 & 10,0 & A resolução concorda com o gabarito do Poliedro nos cinco itens. As semirreações estão corretas nas duas, a resolução em meio ácido e em meio básico. Os valores numéricos também batem: (b) 532 kJ contra 530,4 kJ, (c) 1,35$\times$10$^{-}$$^2$ mol contra 1,36$\times$10$^{-}$$^2$ mol, (d) 1,49$\times$10$^{6}$ kJ contra 1,5 GJ e (e) 837 nm contra 835 nm. Todas as diferenças vêm de o gabarito arredondar N\_A para 6$\times$10$^2$$^3$ e M(H$_{2}$) para 2 g/mol; a resolução usa as constantes fornecidas pela prova. Não há erro de conceito nem de conta em nenhuma das duas. O único ponto a observar é que, no item d, a resolução também cita o limite termodinâmico de 8,00$\times$10$^{5}$ kJ, embora como referência e não como resposta. \\
8 & 10,0 & A resolução está certa nos quatro itens e bate com o gabarito do Poliedro: a) -800 kJ/mol com H2O(g); b) 15 mol de CH4 e 30 mol de O2, ou seja, 240,75 g e 960 g; c) cerca de 0,644 g; d) 200 g. As diferenças em b e c vêm só da massa molar: a resolução usa 16,05 g/mol, calculado com as massas dadas na prova, e o gabarito arredonda para 16 g/mol. O único ponto para levar a um professor é o item a: ali a resolução não escreve a equação multiplicada por 15 nem $\Delta$H = -12 MJ, que o gabarito acrescenta. \\
9 & 10,0 & A resolução bate com o gabarito do Poliedro nos quatro itens. a) Os íons principais são os mesmos, e o CrO$_{4}$$^2$$^{-}$ aparece só no texto, não na lista final. b) Os dois explicam as camadas pela migração iônica: Cu$^2$? (azul) vai para Pt$_{1}$ e Cr$_{2}$O$_{7}$$^2$$^{-}$ (laranja) vai para Pt$_{2}$. c) Pt$_{1}$ é o catodo (2H? + 2e$^{-}$ $\to$ H$_{2}$) e Pt$_{2}$ é o anodo (2H$_{2}$O $\to$ O$_{2}$ + 4H? + 4e$^{-}$). d) P$_{1}$ é negativo, P$_{2}$ é positivo, e a corrente convencional vai de P$_{2}$ a P$_{1}$ pelo circuito externo. Conferi no enunciado original a posição das cores e as ligações P$_{1}$--Pt$_{1}$ e P$_{2}$--Pt$_{2}$, e as duas respostas estão corretas. As diferenças são de detalhe. No item a, a lista final da resolução não inclui o CrO$_{4}$$^2$$^{-}$. Já o gabarito chama de traço o HSO$_{4}$$^{-}$, que na verdade predomina em H$_{2}$SO$_{4}$ diluído. Nenhuma das duas diferenças deve afetar a nota. \\
10 & 10,0 & A resolução está certa e bate com o gabarito do Poliedro nos três itens: a) X = 226, b) Z = 86, c) 3 partículas $\alpha$. Conferi o enunciado na página 12 do PDF e ele corresponde à transcrição. Refiz as contas: o decaimento $\beta$$^{-}$ de $^2$$^3$$^2$$_{8}$$_{9}$Q dá $^2$$^3$$^2$$_{9}$$_{0}$G, com 142 nêutrons. E tem 84 prótons e os mesmos 142 nêutrons, então sua massa é 226, que também é a massa de A. Daí, 238 - 4n = 226 dá n = 3, e 92 - 2$\cdot$3 dá Z = 86. A resolução segue o mesmo caminho do gabarito e ainda verifica a conservação de massa e carga. Não há erro em nenhum dos dois. As únicas ressalvas são de redação: o enunciado diz \textbackslash{}"massa atômica\textbackslash{}" em vez de número de massa e só implicitamente afirma que há apenas emissões $\alpha$. \\
\bottomrule\end{longtable}%
\section*{Questão 1}%
\href{https://poliedroresolve.sistemapoliedro.com.br/vestibulares/ita/2026/ita-2-fase-dia-2-29-10-2025-2-fase-dia-2/dissertativa-1-quimica-2-fase-dia-2-ita-2026}{Resolução do gabarito}\par\smallskip%
\begin{longtable}{@{}p{0.06\textwidth} p{0.33\textwidth} p{0.33\textwidth} p{0.1\textwidth} p{0.08\textwidth}@{}}\toprule Item & Resolução & Gabarito & Veredito & Certo \\ \midrule
a & 0,40 mol/L & 0,4 mol/L & igual & ambos \\
b & 2,0 mol/L & 2,0 mol/L & igual & ambos \\
c & $\approx$ 2,1$\times$10$^{-}$$^2$ mol/L (= 2/95 mol/L) & $\approx$ 2,1$\cdot$10$^{-}$$^2$ mol/L (0,8/38 mol/L) & igual & ambos \\
\bottomrule\end{longtable}%
\begin{itemize}%
\item%
\textbf{a)} Mesmo valor. Conferi: com K\_h = K\_w/K\_A = 10$^{-}$$^2$ e {[}OH$^{-}${]} = {[}HA{]} = 10$^{-}$$^2$, a mistura final tem {[}A$^{-}${]} = 10$^{-}$$^2$ e c\_sal = 2$\times$10$^{-}$$^2$ mol/L. Resolvendo a equação completa x$^2$/(0,02 - x) = 0,01, sai x = 0,0100 exatamente (pH 12,00). Então n = 0,02$\cdot$0,040 = 8$\times$10$^{-}$$^{4}$ mol e C$_{2}$ = 8$\times$10$^{-}$$^{4}$/0,002 = 0,40 mol/L. O gabarito chega ao mesmo valor pela tabela de equilíbrio (5$\cdot$(M\_A)$_{2}$ - 1 = 1). As duas resoluções evitam a aproximação {[}A$^{-}${]} $\approx$ c\_sal, que não vale aqui.%
\item%
\textbf{b)} Diluição de 4 mL para 20 mL, ou seja, fator 5: C$_{1}$ = 5$\cdot$0,40 = 2,0 mol/L.%
\item%
\textbf{c)} 8$\times$10$^{-}$$^{4}$ mol / 0,038 L = 0,02105 mol/L. As frações 2/95 e 0,8/38 são o mesmo número. Coerência: essa base sozinha teria pH $\approx$ 12,32, acima do pH 12 da mistura, como esperado.%
\end{itemize}%
\paragraph{Pontos de dúvida}%
\begin{itemize}%
\item%
\textit{(statement)} O enunciado original escreve ?Solução I? na definição da Solução 2. É erro de digitação para ?Solução 1?, e a resolução declarou essa leitura.%
\item%
\textit{(key)} A transcrição da resolução do gabarito tem defeitos de digitação: ?HA + BOH $\to$ \^{}BA + H$_{2}$O? e ?{[}BA{]} = {[}A$^{-}${]}?, onde deveria ser a concentração analítica do sal, já que no equilíbrio {[}A$^{-}${]} $\neq$ c\_sal. Nada disso muda as contas, que estão certas.%
\item%
\textit{(solution)} Hipóteses implícitas, todas razoáveis e declaradas pela resolução: volumes aditivos (2 + 38 = 40 mL), atividades iguais às concentrações, K\_w = 10$^{-}$$^{1}$$^{4}$ a 25 $^\circ$C e cátion da base forte sem hidrólise.%
\end{itemize}%
\section*{Questão 2}%
\href{https://poliedroresolve.sistemapoliedro.com.br/vestibulares/ita/2026/ita-2-fase-dia-2-29-10-2025-2-fase-dia-2/dissertativa-2-quimica-2-fase-dia-2-ita-2026}{Resolução do gabarito}\par\smallskip%
\begin{longtable}{@{}p{0.06\textwidth} p{0.33\textwidth} p{0.33\textwidth} p{0.1\textwidth} p{0.08\textwidth}@{}}\toprule Item & Resolução & Gabarito & Veredito & Certo \\ \midrule
a & V(O2) $\approx$ 0,61 L (0,6116 L; n(O2) = 0,025 mol, PV = nRT com R = 0,0821 e T = 298 K) & V(O2) $\approx$ 0,6109 L (n(O2) = 0,025 mol, PV = nRT com R = 0,082 e T = 298 K) & equivalente & ambos \\
b & {[}H2O2{]} $\approx$ 0,35 mol$\cdot$L$^{-}$$^{1}$ (0,5 $\times$ e\^{}(-0,3465) = 0,5 $\times$ 0,707 = 0,3535 mol$\cdot$L$^{-}$$^{1}$) & {[}H2O2{]} $\approx$ 0,353 'mol' (é concentração, deveria ser mol/L), calculada como (e\^{}(-0,3465))$^3$ = 0,707$^3$ & equivalente & ambos \\
\bottomrule\end{longtable}%
\begin{itemize}%
\item%
\textbf{a)} Os dois fazem a mesma conta: n(H2O2) = 0,5 $\times$ 0,1 = 0,05 mol, então n(O2) = 0,025 mol, e V = nRT/P. A diferença está só no valor de R: o gabarito usou 0,082 (0,025$\cdot$0,082$\cdot$298 = 0,6109 L) e a resolução usou 0,0821, que é o valor dado na prova (0,025$\cdot$0,0821$\cdot$298 = 0,6116 L). Arredondando, os dois dão 0,61 L. O valor 0,6116 L é um pouco mais fiel às constantes da prova.%
\item%
\textbf{b)} Os dois tratam a reação como de 1ª ordem, por causa da unidade de k em min$^{-}$$^{1}$, e usam kt = 0,0693$\cdot$5 = 0,3465. A resolução faz 0,5$\cdot$0,707 = 0,3535. O gabarito escreve ln{[}H2O2{]} = ln 0,5 - 0,3465 = -0,693 - 0,3465 = -1,0395 = -3$\cdot$0,3465 e calcula 0,707$^3$ = 0,3534. As duas formas valem o mesmo: como 0,707$^2$ $\approx$ 0,5 = e\^{}(-ln 2), temos 0,707$^3$ $\approx$ 0,5$\cdot$0,707. O valor exato, 0,5$\cdot$e\^{}(-0,3465), é 0,35358. Diferença só de arredondamento (0,353 contra 0,3535). O único defeito é do gabarito, que escreveu a unidade como 'mol' em vez de mol$\cdot$L$^{-}$$^{1}$. A resolução está certa nesse ponto.%
\end{itemize}%
\paragraph{Pontos de dúvida}%
\begin{itemize}%
\item%
\textit{(key)} No item b, o gabarito Poliedro escreve o resultado como '0,353 mol', mas o item pede concentração, cuja unidade é mol$\cdot$L$^{-}$$^{1}$. É erro de digitação ou de redação e não muda o valor numérico.%
\item%
\textit{(key)} No item a, o gabarito usa R = 0,082 atm$\cdot$L$\cdot$K$^{-}$$^{1}$$\cdot$mol$^{-}$$^{1}$, mas a prova dá 8,21$\times$10$^{-}$$^2$ (o resultado sai 0,6109 L em vez de 0,6116 L). A diferença não tem importância, porque os dois arredondam para 0,61 L.%
\item%
\textit{(statement)} O enunciado não diz como k foi definida. Se fosse a constante da velocidade da reação, v = -(1/2)d{[}H2O2{]}/dt = k{[}H2O2{]}, o decaimento seria e\^{}(-2kt), e o resultado do item b sairia 0,5$\cdot$e\^{}(-0,693) = 0,25 mol$\cdot$L$^{-}$$^{1}$. Como a prova fornece e\^{}(-0,3465) = 0,707, fica claro que o esperado é kt = 0,3465, ou seja, k é a constante de consumo do H2O2. Gabarito e resolução adotam essa convenção, e a resolução a declara explicitamente.%
\item%
\textit{(solution)} No item a, a resolução supõe decomposição completa, O2 como gás ideal e desconsidera a solubilidade do O2 e a pressão de vapor da água. São hipóteses razoáveis, que o gabarito também adota sem dizer.%
\end{itemize}%
\section*{Questão 3}%
\href{https://poliedroresolve.sistemapoliedro.com.br/vestibulares/ita/2026/ita-2-fase-dia-2-29-10-2025-2-fase-dia-2/dissertativa-3-quimica-2-fase-dia-2-ita-2026}{Resolução do gabarito}\par\smallskip%
\begin{longtable}{@{}p{0.06\textwidth} p{0.33\textwidth} p{0.33\textwidth} p{0.1\textwidth} p{0.08\textwidth}@{}}\toprule Item & Resolução & Gabarito & Veredito & Certo \\ \midrule
a & t $\approx$ 554 s ($\approx$ 5,5$\cdot$10$^2$ s); V no início da precipitação = 1/$\surd$5 $\approx$ 0,446 L, cerca de 553,6 g de água evaporada & $\Delta$t = 1240/2,24 s $\approx$ 553,6 s (V' = 1/2,24 L) & igual & ambos \\
b & {[}A?{]} $\approx$ 6,1$\cdot$10$^{-}$$^3$ mol/L e {[}B$^{-}${]} $\approx$ 1,6$\cdot$10$^{-}$$^3$ mol/L (forma exata: (1+$\surd$3)$\surd$5$\cdot$10$^{-}$$^3$ e ($\surd$3-1)$\surd$5$\cdot$10$^{-}$$^3$; precipitam x $\approx$ 6,35$\cdot$10$^{-}$$^{4}$ mol) & {[}A?{]} = 6,135$\cdot$10$^{-}$$^3$ mol/L e {[}B$^{-}${]} = 1,655$\cdot$10$^{-}$$^3$ mol/L (x = 2,825$\cdot$10$^{-}$$^3$ mol/L) & equivalente & ambos \\
\bottomrule\end{longtable}%
\begin{itemize}%
\item%
\textbf{a)} Resolução e gabarito seguem o mesmo raciocínio. A precipitação começa quando 2$\cdot$10$^{-}$$^{6}$/V$^2$ = 10$^{-}$$^{5}$, ou seja, V = 1/$\surd$5 L. Como a potência e o calor de vaporização têm o mesmo valor numérico (540), t em segundos é igual à massa evaporada em gramas, que é 1000(1 - 1/$\surd$5). Usando $\surd$5 = 2,24, dado da prova, t = 553,6 s. Sem arredondar $\surd$5, t = 552,8 s. A resolução declara $\approx$ 554 s e 553,6 g, o que bate com o gabarito.%
\item%
\textbf{b)} O modelo é o mesmo: V$_{2}$ = V$_{1}$/2, o sistema fica supersaturado (Q = 4K\_PS) e precipita x até {[}A?{]}{[}B$^{-}${]} = K\_PS. A resolução trabalha em mols e o gabarito em concentrações, o que dá no mesmo. Refeito: (2$\cdot$10$^{-}$$^3$ - x)(10$^{-}$$^3$ - x) = 5$\cdot$10$^{-}$$^{7}$, logo x = (3 - $\surd$3)/2$\cdot$10$^{-}$$^3$ mol. Daí {[}A?{]} = 6,109$\cdot$10$^{-}$$^3$ e {[}B$^{-}${]} = 1,637$\cdot$10$^{-}$$^3$ mol/L, e o produto dá exatamente 1,000$\cdot$10$^{-}$$^{5}$. A diferença do gabarito vem de dois arredondamentos. Ele usa $\surd$5 = 2,24 nas concentrações iniciais (8,96 e 4,48$\cdot$10$^{-}$$^3$) e arredonda o termo independente 8,96$\cdot$4,48$\cdot$10$^{-}$$^{6}$ - 10$^{-}$$^{5}$ = 3,014$\cdot$10$^{-}$$^{5}$ para 3$\cdot$10$^{-}$$^{5}$. Com isso, o produto dos seus valores finais é 6,135$\cdot$1,655$\cdot$10$^{-}$$^{6}$ = 1,015$\cdot$10$^{-}$$^{5}$, cerca de 1,5\% acima de K\_PS. Sem o segundo arredondamento, o caminho do gabarito daria 6,115 e 1,635$\cdot$10$^{-}$$^3$. Os valores da resolução estão mais próximos do exato, e as diferenças (0,4\% em A? e 1\% em B$^{-}$) são só de arredondamento.%
\end{itemize}%
\paragraph{Pontos de dúvida}%
\begin{itemize}%
\item%
\textit{(key)} No item b, o gabarito arredonda o termo independente da equação do 2º grau de 3,014$\cdot$10$^{-}$$^{5}$ para 3$\cdot$10$^{-}$$^{5}$. Também usa $\surd$5 $\approx$ 2,24 nas concentrações. Por isso dá {[}B$^{-}${]} = 1,655$\cdot$10$^{-}$$^3$ mol/L, contra 1,637$\cdot$10$^{-}$$^3$ exato, e {[}A?{]}{[}B$^{-}${]} = 1,015$\cdot$10$^{-}$$^{5}$ $\neq$ K\_PS. Com 2 algarismos significativos, o valor do gabarito vira 1,7$\cdot$10$^{-}$$^3$ e o da resolução fica 1,6$\cdot$10$^{-}$$^3$. Essa diferença aparente vem só do arredondamento do gabarito. Convém aceitar valores na faixa de 1,63 a 1,66$\cdot$10$^{-}$$^3$.%
\item%
\textit{(solution)} No item a, há uma pequena inconsistência interna de arredondamento: a resolução escreve 0,554 L = 554 mL evaporados e, logo depois, massa $\approx$ 553,6 g. O valor exato é 552,8 s; com $\surd$5 = 2,24, dado da prova, fica 553,6 s. A resposta '$\approx$ 554 s' é aceitável.%
\item%
\textit{(statement)} Hipóteses implícitas que resolução e gabarito adotam da mesma forma: só a água evapora, o volume do sólido AB é desprezível e o equilíbrio de precipitação é atingido no volume V$_{1}$/2. Pelo enunciado, a 'metade' se refere ao volume do início da precipitação, não ao volume inicial de 1 L, e ambos interpretam assim.%
\end{itemize}%
\section*{Questão 4}%
\href{https://poliedroresolve.sistemapoliedro.com.br/vestibulares/ita/2026/ita-2-fase-dia-2-29-10-2025-2-fase-dia-2/dissertativa-4-quimica-2-fase-dia-2-ita-2026}{Resolução do gabarito}\par\smallskip%
\begin{longtable}{@{}p{0.06\textwidth} p{0.33\textwidth} p{0.33\textwidth} p{0.1\textwidth} p{0.08\textwidth}@{}}\toprule Item & Resolução & Gabarito & Veredito & Certo \\ \midrule
a & m $\approx$ 5,58 g de oxalato amoniacal de nióbio hidratado (1,590 g de Nb2O5 em 10 g de BGN, dividido por 0,285) & m $\approx$ 5,57 g (com n = 9, M = 465 g/mol, M média do BGN = 66,8 g/mol, n(Nb2O5) = 1/167 mol) & equivalente & ambos \\
b & n = 9 (M do sal = 265,82/(2$\cdot$0,285) = 466,4 g/mol; M do sal anidro = 303,0 g/mol; n = 163,4/18,02 = 9,07 $\approx$ 9) & n = 9 (0,285 = 266/(303 + 18n), com 2 mol de oxalato $\to$ 1 mol de Nb2O5) & igual & ambos \\
\bottomrule\end{longtable}%
\begin{itemize}%
\item%
\textbf{a)} Refiz a conta. A massa média do BGN é 0,60$\cdot$60,09 + 0,36$\cdot$56,08 + 0,04$\cdot$265,82 = 66,88 g/mol, e a fração mássica de Nb2O5 é 0,1590, então há 1,590 g de Nb2O5 em 10 g de BGN. A resolução usa direto o resíduo medido de 28,5\%: 1,590/0,285 = 5,579 g. O gabarito usa a fórmula com n = 9 (M = 465,18 g/mol), o que dá exatamente 5,565 g; com os valores arredondados do gabarito (66,8 e 465) sai 5,569 $\approx$ 5,57 g. A diferença entre os dois caminhos é de cerca de 0,25\% e vem só de usar o teor medido em vez da fórmula com n inteiro. Os dois caminhos estão certos e os resultados se equivalem. A própria resolução faz a verificação com n = 9 e obtém 5,56 g.%
\item%
\textbf{b)} Refiz a conta. M(NH4{[}NbO(C2O4)2{]}) = 303,00 g/mol e 0,285 = 265,82/{[}2(303,00 + 18,02n){]}, o que dá n = 9,06 $\approx$ 9. Com n = 9, o teor teórico de Nb2O5 é 28,57\%; com n = 8 seria 29,7\% e com n = 10, 27,5\%. Logo, n = 9 é o único valor inteiro compatível com os 28,5\%. A equação escrita no gabarito omite o fator 2 do denominador. Tomada ao pé da letra, ela daria n $\approx$ 35. Mesmo assim, o gabarito chega a n $\approx$ 9 e M = 465 g/mol, valores coerentes com o fator 2 que ele mesmo declara (2 mol de oxalato $\to$ 1 mol de Nb2O5). Trata{-}se, portanto, de erro de digitação, e o resultado está correto.%
\end{itemize}%
\paragraph{Pontos de dúvida}%
\begin{itemize}%
\item%
\textit{(key)} A equação do gabarito no item b está escrita como 0,285 = 266/(303 + 18n), sem o fator 2 no denominador. Ao pé da letra, ela dá n $\approx$ 35. O texto do gabarito (2 mol de oxalato $\to$ 1 mol de Nb2O5) e o resultado n $\approx$ 9 mostram que é só erro de digitação, e a forma correta é 0,285 = 266/{[}2(303 + 18n){]}. A resposta final não muda.%
\item%
\textit{(statement)} No item a, o enunciado não diz se a massa de sal deve vir do teor medido (28,5\%, que dá 5,58 g) ou da fórmula com n inteiro (n = 9, que dá 5,56--5,57 g). As duas leituras são aceitáveis e diferem em menos de 0,5\%.%
\item%
\textit{(key)} Em atmosfera oxidante a 1000 $^\circ$C, o gabarito lista NH3 e CO entre os gases liberados. Em oxidação completa, o esperado seria CO2, H2O e N2/NOx, que é o que a resolução descreve. Isso não afeta o resultado numérico.%
\end{itemize}%
\section*{Questão 5}%
\href{https://poliedroresolve.sistemapoliedro.com.br/vestibulares/ita/2026/ita-2-fase-dia-2-29-10-2025-2-fase-dia-2/dissertativa-5-quimica-2-fase-dia-2-ita-2026}{Resolução do gabarito}\par\smallskip%
\begin{longtable}{@{}p{0.06\textwidth} p{0.33\textwidth} p{0.33\textwidth} p{0.1\textwidth} p{0.08\textwidth}@{}}\toprule Item & Resolução & Gabarito & Veredito & Certo \\ \midrule
a & 3 arranjos lineares, um para cada átomo central: SCN$^{-}$ (C central), CNS$^{-}$ (N central) e CSN$^{-}$ (S central); só o SCN$^{-}$ é estável. O anel triangular foi descartado. & 3 arranjos, conforme o átomo central: {[}CNS{]}$^{-}$, {[}NCS{]}$^{-}$ ou {[}CSN{]}$^{-}$ (na notação do enunciado: {[}XGQ{]}$^{-}$, {[}GXQ{]}$^{-}$ ou {[}XQG{]}$^{-}$). & igual & ambos \\
b & Trigonal planar: BF$_{3}$. Pirâmide trigonal: NH$_{3}$. Cita também SO$_{3}$, CH$_{2}$O e NF$_{3}$ como outros exemplos. & Trigonal planar: BF$_{3}$, BH$_{3}$ ou SO$_{3}$ (ZY$_{3}$, ZA$_{3}$ ou QR$_{3}$). Pirâmide trigonal: NH$_{3}$ ou NF$_{3}$ (GA$_{3}$ ou GY$_{3}$). & igual & ambos \\
c & O composto é o fluorobenzeno, que sofre nitração por substituição eletrofílica aromática. Produtos: 1{-}fluoro{-}4{-}nitrobenzeno (para, majoritário, $\approx$85--90\%) e 1{-}fluoro{-}2{-}nitrobenzeno (orto, minoritário), além de H$_{2}$O. As estruturas foram desenhadas corretamente: NO$_{2}$ oposto ao F no para e vizinho ao F no orto. & Fluorobenzeno + HNO$_{3}$/H$_{2}$SO$_{4}$ dá 1{-}fluoro{-}2{-}nitrobenzeno (orto) e 1{-}fluoro{-}4{-}nitrobenzeno (para), porque o F orienta para orto/para. & igual & ambos \\
d & (ii) H$_{2}$S < (iii) CH$_{3}$OH < (i) H$_{2}$O & H$_{2}$S < CH$_{3}$OH < H$_{2}$O, ou seja, (ii) < (iii) < (i) & igual & ambos \\
\bottomrule\end{longtable}%
\begin{itemize}%
\item%
\textbf{a)} Os dois contam 3 arranjos, um para cada átomo central, e não contam o anel. Conferi as cargas formais que a resolução usa como justificativa extra. Com C central, S=C=N$^{-}$ dá (0,0,-1). Com N central, C$^{-}$$\equiv$N?--S$^{-}$ dá (-1,+1,-1). Com S central, C=S=N$^{-}$ dá (-2,+2,-1); as outras estruturas com octeto completo, C$\equiv$S--N e C--S$\equiv$N, também somam |CF| = 5. Essas análises estão certas, mas o item não as pede.%
\item%
\textbf{b)} BF$_{3}$ e NH$_{3}$ estão entre as respostas que o gabarito aceita. Pela VSEPR, o BF$_{3}$ tem 3 domínios e nenhum par isolado, por isso é trigonal planar. O NH$_{3}$ tem 3 ligações e 1 par isolado, por isso é piramidal. Os exemplos extras da resolução também valem. O SO$_{3}$ está na lista do gabarito. O CH$_{2}$O é trigonal planar, embora não equilátero.%
\item%
\textbf{c)} Os dois apresentam os mesmos produtos principais, orto e para. A resolução acrescenta a proporção para > orto, que bate com os dados experimentais: cerca de 87--91\% para, 9--13\% orto e meta abaixo de 1\%. Acrescenta também a água como coproduto. Nenhum dos dois acréscimos contradiz o gabarito. O desenho em TikZ coloca o F a 90$^\circ$ e o NO$_{2}$ a -90$^\circ$ (para) ou a 30$^\circ$ (orto), o que está correto.%
\item%
\textbf{d)} A ordem e a justificativa são as mesmas nos dois: o H$_{2}$S tem só dipolo{-}dipolo e dispersão, e a H$_{2}$O faz mais ligações de hidrogênio por molécula que o metanol. Os dados citados pela resolução estão corretos. Pontos de ebulição: H$_{2}$S $\approx$ -60 $^\circ$C, CH$_{3}$OH $\approx$ 65 $^\circ$C, H$_{2}$O = 100 $^\circ$C. $\Delta$Hvap: $\approx$ 18,7, 35 e 40,7 kJ/mol.%
\end{itemize}%
\paragraph{Pontos de dúvida}%
\begin{itemize}%
\item%
\textit{(statement)} No item a, a expressão 'arranjos atômicos' é ambígua. O gabarito e a resolução contam só as cadeias abertas (3). Se o arranjo cíclico C--N--S fosse aceito como arranjo distinto, a resposta seria 4. Os dois descartam o anel, e a resolução justifica isso: tensão angular e geometria linear esperada para um sistema de 16 elétrons de valência.%
\item%
\textit{(key)} No item b, o gabarito aceita BH$_{3}$ como exemplo trigonal planar. O BH$_{3}$ isolado é instável e se dimeriza em B$_{2}$H$_{6}$, mas como espécie molecular em VSEPR a resposta é aceitável. Isso não afeta a resolução, que deu BF$_{3}$.%
\item%
\textit{(solution)} No item c, a resolução chama o isômero orto de 'minoritário', enquanto o gabarito lista orto e para juntos como produtos principais. Como a resolução desenha os dois e o enunciado diz 'produto(s) principal(is)', não deve haver desconto.%
\end{itemize}%
\section*{Questão 6}%
\href{https://poliedroresolve.sistemapoliedro.com.br/vestibulares/ita/2026/ita-2-fase-dia-2-29-10-2025-2-fase-dia-2/dissertativa-6-quimica-2-fase-dia-2-ita-2026}{Resolução do gabarito}\par\smallskip%
\begin{longtable}{@{}p{0.06\textwidth} p{0.33\textwidth} p{0.33\textwidth} p{0.1\textwidth} p{0.08\textwidth}@{}}\toprule Item & Resolução & Gabarito & Veredito & Certo \\ \midrule
a & pKa(H3O+) = -1,70; pKb(CH3CC-) = -11,00; K1 = 10\^{}(-26,70) $\approx$ 2$\times$10\^{}(-27) < 1: não favorece os produtos & pKa(H3O+) = -1,70; pKb(CH3CC-) = -11; Keq = 10\^{}(-1,70-25) < 1: não favorece os produtos (deslocado para a esquerda) & igual & ambos \\
b & pKa(NH4+) = 9,25; K2 = 10\^{}(-15,75) $\approx$ 1,8$\times$10\^{}(-16) < 1: não favorece os produtos & pKa(NH4+) = 14 - 4,75 = 9,25; Keq = 10\^{}(9,25-25) < 1: não favorece os produtos (deslocado para a esquerda) & igual & ambos \\
c & pKa(NH3) = 36,00; K3 = 10\^{}11 > 1: favorece os produtos & pKa(NH3) = 14 - (-22) = 36; Keq = 10\^{}(36-25) = 10\^{}11 > 1: favorece os produtos (deslocado para a direita) & igual & ambos \\
\bottomrule\end{longtable}%
\begin{itemize}%
\item%
\textbf{a)} Os valores são os mesmos: pKa(H3O+) = 14 - 15,70 = -1,70 e log K = -1,70 - 25,00 = -26,70. A conclusão também é a mesma. Conferi 10\^{}(-26,7) $\approx$ 2,0$\times$10\^{}(-27). O gabarito comenta que o pKb da água deveria ser 14, mas isso não muda o resultado.%
\item%
\textbf{b)} Os valores são idênticos: log K = 9,25 - 25 = -15,75. Conferi 10\^{}(-15,75) $\approx$ 1,78$\times$10\^{}(-16). A conclusão é a mesma.%
\item%
\textbf{c)} Os valores são idênticos. A resolução tira o Na+ espectador e escreve a reação efetiva CH3CCH + NH2- ? CH3CC- + NH3, a mesma que o gabarito usa. Ela ainda confere o resultado com os pKb: 14 - 25 + 22 = 11.%
\end{itemize}%
\paragraph{Pontos de dúvida}%
\begin{itemize}%
\item%
\textit{(key)} O gabarito diz que o pKb da água deveria ser 14, citando J. Chem. Educ. 2017. O valor 15,70 do enunciado segue a convenção que inclui {[}H2O{]} = 55,5 mol/L, que dá pKa(H3O+) $\approx$ -1,74. É uma questão de convenção. Com o valor 15,70 do enunciado, que a resolução e o gabarito usam, pKa(H3O+) = -1,70. Com 14, daria pKa(H3O+) = 0. Em nenhum dos casos a conclusão do item a muda.%
\item%
\textit{(statement)} O enunciado manda usar pKa + pKb = 14, que vale para a escala aquosa a 25 $^\circ$C, também para espécies que não existem em água: NH2- e o acetileto, com pKb negativo, e o NH3 atuando como ácido. Esses valores servem só para comparar forças relativas. A resolução declara isso e observa que, na prática, o NaNH2 é usado em NH3 líquida. O gabarito adota a mesma hipótese sem comentar.%
\item%
\textit{(solution)} 'Demais espécies' admite mais de uma leitura. A resolução, como o gabarito, calcula só os quatro valores conjugados que entram nas reações. Não calcula, por exemplo, o pKa da água atuando como ácido (par H2O/OH-). Não parece cobrado, já que o gabarito apresenta exatamente os mesmos quatro valores.%
\end{itemize}%
\section*{Questão 7}%
\href{https://poliedroresolve.sistemapoliedro.com.br/vestibulares/ita/2026/ita-2-fase-dia-2-29-10-2025-2-fase-dia-2/dissertativa-7-quimica-2-fase-dia-2-ita-2026}{Resolução do gabarito}\par\smallskip%
\begin{longtable}{@{}p{0.06\textwidth} p{0.33\textwidth} p{0.33\textwidth} p{0.1\textwidth} p{0.08\textwidth}@{}}\toprule Item & Resolução & Gabarito & Veredito & Certo \\ \midrule
a & Meio ácido: cátodo 4H? + 4e$^{-}$ $\to$ 2H$_{2}$; ânodo 2H$_{2}$O $\to$ O$_{2}$ + 4H? + 4e$^{-}$ (também dá a forma em meio básico: cátodo 4H$_{2}$O + 4e$^{-}$ $\to$ 2H$_{2}$ + 4OH$^{-}$; ânodo 4OH$^{-}$ $\to$ O$_{2}$ + 2H$_{2}$O + 4e$^{-}$) & Cátodo: 2H$_{2}$O + 2e$^{-}$ $\to$ H$_{2}$ + 2OH$^{-}$; Ânodo: 2H$_{2}$O $\to$ O$_{2}$ + 4H? + 4e$^{-}$ & equivalente & ambos \\
b & $\approx$ 532 kJ por mol de H$_{2}$ (2 mol de fótons de 450 nm; E\_f = 4,42$\times$10$^{-}$$^{1}$$^{9}$ J) & $\approx$ 530,4 kJ (5,304$\times$10$^{5}$ J) por mol de H$_{2}$ & equivalente & ambos \\
c & $\approx$ 1,35$\times$10$^{-}$$^2$ mol de H$_{2}$ & $\approx$ 1,36$\times$10$^{-}$$^2$ mol de H$_{2}$ & equivalente & ambos \\
d & $\approx$ 1,49$\times$10$^{6}$ kJ (n = 5650/2,02 = 2797 mol $\times$ 532 kJ/mol); cita como referência o limite termodinâmico de 8,00$\times$10$^{5}$ kJ & n = 5650/2 = 2825 mol; E = 530,4$\cdot$2825 = 1 498 380 kJ $\approx$ 1,5 GJ & equivalente & ambos \\
e & $\lambda$\_máx $\approx$ 837 nm & $\lambda$ $\approx$ 835 nm & equivalente & ambos \\
\bottomrule\end{longtable}%
\begin{itemize}%
\item%
\textbf{a)} As duas respostas trazem semirreações corretas de redução da água/H? a H$_{2}$ no cátodo e de oxidação a O$_{2}$ no ânodo. A resolução apresenta as formas ácida e básica, cada uma coerente consigo mesma. O gabarito mistura cátodo em meio básico com ânodo em meio ácido. Mesmo assim, a soma (com o cátodo $\times$2) dá 2H$_{2}$O $\to$ 2H$_{2}$ + O$_{2}$ depois de H? + OH$^{-}$ $\to$ H$_{2}$O. A forma do cátodo da resolução para o meio básico é exatamente a do gabarito multiplicada por 2.%
\item%
\textbf{b)} O método é o mesmo nas duas: E = 2$\cdot$N\_A$\cdot$hc/$\lambda$. A diferença vem só do valor de N\_A. Com 6,02$\times$10$^2$$^3$, que é o valor dado na prova, sai 532,2 kJ; com 6$\times$10$^2$$^3$, que o gabarito arredondou, sai 530,4 kJ. A resolução usa a constante fornecida e por isso é um pouco mais fiel aos dados.%
\item%
\textbf{c)} Nas duas, E = 10$\cdot$0,2$\cdot$3600 = 7200 J e n = 7200/E(b). A conta dá 7200/532168 = 1,353$\times$10$^{-}$$^2$, contra 7200/530400 = 1,357$\times$10$^{-}$$^2$ no gabarito. A diferença é só de arredondamento, herdada do N\_A usado no item b.%
\item%
\textbf{d)} Na leitura adotada a resolução usa a energia dos fótons de 450 nm por mol de H$_{2}$, a mesma do gabarito. Recalculando: 2797$\cdot$532,2 $\approx$ 1,488$\times$10$^{6}$ kJ, enquanto 2825$\cdot$530,4 = 1,498$\times$10$^{6}$ kJ. Ambos ficam em $\approx$ 1,5$\times$10$^{6}$ kJ = 1,5 GJ. O gabarito usa M(H$_{2}$) = 2 g/mol, mas pelos dados da prova seria 2,02 g/mol. A resolução também cita o valor termodinâmico de 8,00$\times$10$^{5}$ kJ, porém só como referência; a resposta em destaque é a mesma do gabarito.%
\item%
\textbf{e)} A fórmula é a mesma: $\lambda$ = 2$\cdot$N\_A$\cdot$h$\cdot$c/(286$\times$10$^3$ J). Com N\_A = 6,02$\times$10$^2$$^3$ dá 837,3 nm; com 6$\times$10$^2$$^3$, como no gabarito, dá 834,5 nm. A diferença de cerca de 0,3\% vem só da constante. As duas respostas consideram a combustão formando H$_{2}$O(?), com 286 kJ/mol.%
\end{itemize}%
\paragraph{Pontos de dúvida}%
\begin{itemize}%
\item%
\textit{(key)} O gabarito arredonda N\_A para 6$\times$10$^2$$^3$ em vez de usar 6,02$\times$10$^2$$^3$, que é o valor dado na prova, e usa M(H$_{2}$) = 2 g/mol em vez de 2,02 g/mol (H = 1,01). Por isso os números dele saem um pouco diferentes: 530,4 kJ contra 532 kJ, 1,36$\times$10$^{-}$$^2$ contra 1,35$\times$10$^{-}$$^2$, 2825 mol contra 2797 mol e 835 nm contra 837 nm. Nenhuma dessas diferenças muda a resposta.%
\item%
\textit{(key)} No item a, o gabarito escreve o cátodo em meio básico (com OH$^{-}$) e o ânodo em meio ácido (com H?), o que é incoerente quanto ao meio. Além disso, deixa o número de elétrons diferente nas duas semirreações (2e$^{-}$ e 4e$^{-}$). Como o enunciado não fixa o meio, isso não chega a ser erro, mas o mais rigoroso seria usar o mesmo meio nos dois eletrodos.%
\item%
\textit{(solution)} No item d, a resolução fica com 1,49$\times$10$^{6}$ kJ, que coincide com o gabarito, mas também apresenta 8,00$\times$10$^{5}$ kJ como 'limite termodinâmico' para quem lesse 'mínima' de outro modo. A resposta principal está claramente marcada, então o risco é pequeno. Ainda assim, um corretor rigoroso pode ver isso como duas respostas para o mesmo item.%
\item%
\textit{(statement)} O enunciado não diz se a água é formada como líquido ou como vapor na combustão do item e. As duas respostas usam H$_{2}$O(?) e 286 kJ/mol, de acordo com o $\Delta$H$^\circ$ fornecido. Com H$_{2}$O(g), cerca de 242 kJ/mol, o comprimento de onda daria perto de 990 nm, mas esse dado não aparece na prova.%
\end{itemize}%
\section*{Questão 8}%
\href{https://poliedroresolve.sistemapoliedro.com.br/vestibulares/ita/2026/ita-2-fase-dia-2-29-10-2025-2-fase-dia-2/dissertativa-8-quimica-2-fase-dia-2-ita-2026}{Resolução do gabarito}\par\smallskip%
\begin{longtable}{@{}p{0.06\textwidth} p{0.33\textwidth} p{0.33\textwidth} p{0.1\textwidth} p{0.08\textwidth}@{}}\toprule Item & Resolução & Gabarito & Veredito & Certo \\ \midrule
a & CH4(g) + 2 O2(g) $\to$ CO2(g) + 2 H2O(g); $\Delta$Hf(H2O,g) = -286 + 45 = -241 kJ/mol; $\Delta$Hc = -800 kJ/mol de CH4 & CH4(g) + 2 O2(g) $\to$ CO2(g) + 2 H2O(g), $\Delta$H = -800 kJ/mol de CH4 ($\Delta$Hf H2O(g) = -241). Para 12 MJ: 15 CH4 + 30 O2 $\to$ 15 CO2 + 30 H2O(g), $\Delta$H = -12 000 kJ & equivalente & ambos \\
b & m(CH4) = 15$\cdot$16,05 = 240,75 g; m(O2) = 30$\cdot$32,00 = 960 g & m(CH4) = 15$\cdot$16 = 240 g; m(O2) = 960 g & equivalente & ambos \\
c & E útil = 40$\cdot$240,75 = 9630 J; Q pré = 9630/0,30 = 32,1 kJ; n = 32,1/800 = 0,040125 mol; m $\approx$ 0,644 g de CH4(g) & Q útil = 240$\cdot$40 = 9,6 kJ; Q comb = 32 kJ; m = 32/800$\cdot$16 = 0,64 g & equivalente & ambos \\
d & Q = 0,0025$\cdot$12 MJ = 3,0$\cdot$10\^{}4 J; m = 3,0$\cdot$10\^{}4/((5/3)$\cdot$90) = 200 g de CH4(l) & Q = 30 kJ = m$\cdot$(5/3)$\cdot$90 $\to$ m = 200 g (mínima por usar T máx = 188 K) & igual & ambos \\
\bottomrule\end{longtable}%
\begin{itemize}%
\item%
\textbf{a)} A equação, o estado gasoso da água e $\Delta$Hc = {[}-393 + 2(-241){]} - (-75) = -800 kJ/mol coincidem. O gabarito também escreve a equação multiplicada por 15, com $\Delta$H = -12 MJ. A resolução não faz isso no item a, mas calcula os 15 mol logo no item b. O conteúdo químico é o mesmo.%
\item%
\textbf{b)} Ambos usam n(CH4) = 12000/800 = 15 mol e n(O2) = 30 mol. A diferença de 0,75 g vem só da massa molar: a resolução usa M(CH4) = 12,01 + 4$\cdot$1,01 = 16,05 g/mol, que são as massas dadas na prova, e o gabarito arredonda para 16. É diferença de arredondamento, e o valor da resolução é até mais fiel aos dados. O O2 dá 960 g nos dois.%
\item%
\textbf{c)} O raciocínio é o mesmo: a energia útil é 40 J/g vezes a massa de metano da câmara principal, o calor da pré{-}câmara é a energia útil dividida por 0,30, e o metano queimado libera 800 kJ/mol. Recalculei com 16,05 g/mol: 9630 J $\to$ 32 100 J $\to$ 0,040125 mol $\to$ 0,6440 g. Com 16 g/mol, como no gabarito, sai 0,64 g. A diferença é só de arredondamento.%
\item%
\textbf{d)} Recalculei: 0,0025$\cdot$12$\cdot$10\^{}6 J = 30 000 J e (5/3)$\cdot$(188 - 98) = 150 J/g. Então m = 200 g. A justificativa da massa mínima (aquecer até a temperatura máxima de 188 K) é a mesma nos dois.%
\end{itemize}%
\paragraph{Pontos de dúvida}%
\begin{itemize}%
\item%
\textit{(statement)} O item a pede para "calcular a entalpia da combustão, conforme o requisito I". Isso dá margem a duas leituras: $\Delta$H por mol (-800 kJ/mol, com H2O gasosa) ou $\Delta$H do sistema que atende a 12 MJ (-12 000 kJ, com a equação multiplicada por 15). O gabarito traz as duas. A resolução só traz o valor por mol no item a e chega aos 15 mol no item b. Um corretor rigoroso poderia descontar uma fração no item a por falta da equação escalonada e de $\Delta$H = -12 MJ.%
\item%
\textit{(key)} O gabarito usa M(CH4) = 16 g/mol, quando as massas dadas na prova dão 16,05 g/mol. Por isso tem 240 g contra 240,75 g no item b e 0,64 g contra 0,644 g no item c. É só arredondamento e não muda o critério.%
\item%
\textit{(statement)} Os dois adotam as mesmas hipóteses implícitas no item c, que são razoáveis. (i) A turbobomba pressuriza só o metano da câmara principal (240 g), sem contar o metano da pré{-}câmara. (ii) O calor da pré{-}câmara é separado dos 12 MJ da câmara principal, embora os gases da pré{-}câmara entrem nela. (iii) Na pré{-}câmara a água também sai gasosa, com 800 kJ/mol. (iv) O O2 consumido na pré{-}câmara (cerca de 0,08 mol) não entra nos 960 g do item b.%
\item%
\textit{(statement)} No original, a unidade aparece como "KJ$\cdot$mol$^{-}$$^{1}$", com K maiúsculo. É erro tipográfico sem efeito na resolução. A transcrição confere com o PDF (página 10).%
\end{itemize}%
\section*{Questão 9}%
\href{https://poliedroresolve.sistemapoliedro.com.br/vestibulares/ita/2026/ita-2-fase-dia-2-29-10-2025-2-fase-dia-2/dissertativa-9-quimica-2-fase-dia-2-ita-2026}{Resolução do gabarito}\par\smallskip%
\begin{longtable}{@{}p{0.06\textwidth} p{0.33\textwidth} p{0.33\textwidth} p{0.1\textwidth} p{0.08\textwidth}@{}}\toprule Item & Resolução & Gabarito & Veredito & Certo \\ \midrule
a & Cu$^2$?, K?, H?, Cr$_{2}$O$_{7}$$^2$$^{-}$, SO$_{4}$$^2$$^{-}$, HSO$_{4}$$^{-}$ e OH$^{-}$ (traços). No texto, cita também CrO$_{4}$$^2$$^{-}$ e HCrO$_{4}$$^{-}$ como traços, mas não os põe na lista final. & Principais: H? (H$_{3}$O?), SO$_{4}$$^2$$^{-}$, Cu$^2$?, K?, Cr$_{2}$O$_{7}$$^2$$^{-}$. Em quantidades muito pequenas: OH$^{-}$, HSO$_{4}$$^{-}$, CrO$_{4}$$^2$$^{-}$. & equivalente & ambos \\
b & Migração iônica sob o campo elétrico. O Cu$^2$? (azul) sobe pelo braço de Pt$_{1}$ (catodo) e o Cr$_{2}$O$_{7}$$^2$$^{-}$ se afasta dele, formando a camada azul. O Cr$_{2}$O$_{7}$$^2$$^{-}$ (laranja) sobe pelo braço de Pt$_{2}$ (anodo) e o Cu$^2$? desce, formando a camada alaranjada. No fundo do U os íons ainda não se separaram, e a cor continua marrom. & Migração iônica: Pt$_{1}$ (polo -) atrai os cátions K? e H$_{3}$O? (incolores) e Cu$^2$? (azul), e a solução perto de Pt$_{1}$ fica azul. Pt$_{2}$ (polo +) atrai os ânions SO$_{4}$$^2$$^{-}$ (incolor) e Cr$_{2}$O$_{7}$$^2$$^{-}$ (alaranjado), e a solução perto de Pt$_{2}$ fica alaranjada. & igual & ambos \\
c & Pt$_{1}$ é o catodo: 2H? + 2e$^{-}$ $\to$ H$_{2}$(g). Pt$_{2}$ é o anodo: 2H$_{2}$O $\to$ O$_{2}$(g) + 4H? + 4e$^{-}$. Os eletrodos estão imersos no H$_{2}$SO$_{4}$. Se o Cu$^2$? chegasse ao catodo, depositaria Cu(s). & Pt$_{1}$ é o catodo (polo -, redução): 2H?(aq) + 2e$^{-}$ $\to$ H$_{2}$(g). Pt$_{2}$ é o anodo (polo +, oxidação): 2H$_{2}$O(?) $\to$ O$_{2}$(g) + 4H?(aq) + 4e$^{-}$. As placas ficam em contato com a solução de H$_{2}$SO$_{4}$. & igual & ambos \\
d & P$_{1}$ é negativo e P$_{2}$ é positivo. Elétrons: P$_{1}$ $\to$ Pt$_{1}$ e Pt$_{2}$ $\to$ P$_{2}$. Corrente convencional: P$_{2}$ (+) $\to$ Pt$_{2}$ $\to$ solução $\to$ Pt$_{1}$ $\to$ P$_{1}$ (-). & P$_{1}$ é negativo e P$_{2}$ é positivo. Os elétrons vão de P$_{1}$ para P$_{2}$ (sentido real). A corrente convencional vai de P$_{2}$ para P$_{1}$ pelo circuito externo. & equivalente & ambos \\
\bottomrule\end{longtable}%
\begin{itemize}%
\item%
\textbf{a)} Os íons principais são os mesmos. A única diferença é que a lista final da resolução deixa de fora o CrO$_{4}$$^2$$^{-}$, que o gabarito põe entre os traços. A resolução, porém, discute o equilíbrio Cr$_{2}$O$_{7}$$^2$$^{-}$/CrO$_{4}$$^2$$^{-}$ e diz que o CrO$_{4}$$^2$$^{-}$ (e o HCrO$_{4}$$^{-}$) existe só como traço, então o conteúdo é o mesmo. A resolução trata o HSO$_{4}$$^{-}$ como espécie relevante e o gabarito o põe entre os traços. Do ponto de vista químico, a resolução está mais correta: com Ka2 $\approx$ 1$\times$10$^{-}$$^2$, o HSO$_{4}$$^{-}$ é cerca de 59\% do sulfato total em H$_{2}$SO$_{4}$ 0,01 mol/L e cerca de 92\% em 0,1 mol/L. Isso não muda a resposta do item.%
\item%
\textbf{b)} A explicação é a mesma: cada íon colorido migra para o eletrodo de carga oposta, e isso separa o azul (Cu$^2$?) do alaranjado (Cr$_{2}$O$_{7}$$^2$$^{-}$). Pela Figura (B) do original, o azul fica sob Pt$_{1}$ (braço esquerdo) e o alaranjado sob Pt$_{2}$ (braço direito), como as duas dizem. A resolução detalha mais, porque explica também por que cada camada fica 'pura': o íon de carga oposta se afasta. Mesma conclusão.%
\item%
\textbf{c)} A atribuição dos eletrodos e as duas semirreações são as mesmas, e a justificativa também: as placas estão na camada incolor de H$_{2}$SO$_{4}$. Os potenciais que a resolução cita estão corretos: Cu$^2$?/Cu +0,34 V, Cr$_{2}$O$_{7}$$^2$$^{-}$ +1,33 V, O$_{2}$/H$_{2}$O +1,23 V e S$_{2}$O$_{8}$$^2$$^{-}$/SO$_{4}$$^2$$^{-}$ $\approx$ +2,01 V. A observação sobre o depósito de Cu é um comentário a mais e não contradiz o gabarito.%
\item%
\textbf{d)} Os sinais dos polos são os mesmos. O sentido da corrente também é o mesmo: a resolução descreve o caminho completo (P$_{2}$ $\to$ Pt$_{2}$ $\to$ solução $\to$ Pt$_{1}$ $\to$ P$_{1}$), e o gabarito resume como 'de P$_{2}$ para P$_{1}$'. Os dois indicam também o sentido real dos elétrons, oposto ao da corrente convencional. No original, P$_{1}$ está ligado a Pt$_{1}$ (esquerda) e P$_{2}$ a Pt$_{2}$ (direita), o que confirma as conclusões.%
\end{itemize}%
\paragraph{Pontos de dúvida}%
\begin{itemize}%
\item%
\textit{(key)} O gabarito põe o HSO$_{4}$$^{-}$ entre as espécies 'em quantidades muito pequenas'. Com Ka2(H$_{2}$SO$_{4}$) $\approx$ 1$\times$10$^{-}$$^2$, o HSO$_{4}$$^{-}$ é a forma predominante do sulfato em H$_{2}$SO$_{4}$ diluído: cerca de 59\% em 0,01 mol/L e cerca de 92\% em 0,1 mol/L. A classificação do gabarito é discutível, mas não muda o que se pede no item a.%
\item%
\textit{(solution)} A lista final do item a deixa de fora o CrO$_{4}$$^2$$^{-}$ (e o HCrO$_{4}$$^{-}$), que o gabarito põe entre os traços. A resolução cita essas espécies no texto. Como o enunciado pede 'todas as espécies iônicas', um corretor rigoroso poderia exigir que aparecessem na lista final.%
\item%
\textit{(solution)} No item b, a resolução diz que as fronteiras ficam nítidas também porque as soluções têm 'densidades diferentes'. Essa hipótese não está no enunciado nem é necessária. A explicação que se espera é a migração iônica, e ela está presente.%
\item%
\textit{(statement)} O item c pede a reação que ocorre 'preferencialmente' no catodo. Pelos potenciais, o Cu$^2$? (+0,34 V) seria reduzido antes do H? se estivesse em contato com a placa. O gabarito e a resolução escolhem 2H? + 2e$^{-}$ $\to$ H$_{2}$ porque, na Figura (B), Pt$_{1}$ está na camada incolor de H$_{2}$SO$_{4}$. Um aluno que escrevesse Cu$^2$? + 2e$^{-}$ $\to$ Cu poderia gerar discussão na correção.%
\end{itemize}%
\section*{Questão 10}%
\href{https://poliedroresolve.sistemapoliedro.com.br/vestibulares/ita/2026/ita-2-fase-dia-2-29-10-2025-2-fase-dia-2/dissertativa-10-quimica-2-fase-dia-2-ita-2026}{Resolução do gabarito}\par\smallskip%
\begin{longtable}{@{}p{0.06\textwidth} p{0.33\textwidth} p{0.33\textwidth} p{0.1\textwidth} p{0.08\textwidth}@{}}\toprule Item & Resolução & Gabarito & Veredito & Certo \\ \midrule
a & X = 226 & X = 226 & igual & ambos \\
b & Z = 86 & Z = 86 & igual & ambos \\
c & 3 partículas alfa & 3 partículas alfa & igual & ambos \\
\bottomrule\end{longtable}%
\begin{itemize}%
\item%
\textbf{a)} Confere. O decaimento $\beta$$^{-}$ de $^2$$^3$$^2$$_{8}$$_{9}$Q dá $^2$$^3$$^2$$_{9}$$_{0}$G, com N = 142. E é isótono de G e tem Z = 84, então A\_E = 84 + 142 = 226. A é isóbaro de E, logo X = 226. A resolução e o gabarito seguem o mesmo raciocínio.%
\item%
\textbf{b)} Confere. Pela conservação da carga, 92 = 2$\cdot$3 + Z, logo Z = 86. A resolução também confere a massa (238 = 226 + 12) e a carga, e identifica o núcleo como $^2$$^2$$^{6}$Rn.%
\item%
\textbf{c)} Confere. Pela conservação da massa, 238 = 4n + 226, logo n = 3, um valor inteiro e coerente com a hipótese de haver só emissões $\alpha$.%
\end{itemize}%
\paragraph{Pontos de dúvida}%
\begin{itemize}%
\item%
\textit{(statement)} O enunciado usa "massa atômica igual a 238" quando o certo seria número de massa. Também não diz explicitamente que J emite apenas partículas $\alpha$, mas fica implícito em "sofre decaimentos de partículas alfa". A resolução e o gabarito assumem que só há emissões $\alpha$, e é isso que leva a Z = 86. Se houvesse $\beta$$^{-}$ na cadeia, Z não ficaria determinado. Na série real do U{-}238, o núcleo com massa 226 é o $^2$$^2$$^{6}$Ra (Z = 88), obtido com 3$\alpha$ + 2$\beta$$^{-}$, e não o $^2$$^2$$^{6}$Rn. Como os elementos são fictícios e o enunciado fala só de $\alpha$, a leitura adotada é a pretendida.%
\item%
\textit{(solution)} A resolução usa a letra A ao mesmo tempo para o número de massa (A = Z + N) e para o nome do elemento A. Isso pode confundir na leitura, mas não muda nenhum resultado.%
\end{itemize}%
\end{document}